% Imporditud: tools/import_eksperimendid.py
% Allikas: Eesti füüsikaolümpiaadi eksperimendi ülesannete
%   kogu (Rael Kalda, 2023), ülesanne Ü78, lahendus L78.
\setAuthor{EFO žürii}
\setRound{lõppvoor}
\setYear{1998}
\setNumber{G E2}
\setDifficulty{3}
\setTopic{vedelike mehaanika}
\setVahendid{3 numbriga anumat vedelikega, pulgake, tükike plastiliini, tualettpaberit kuivatamiseks}
\setPunktid{10}
\setAllikas{Eksperimendi kogu Ü78, lahendus lk 67}
\setLahendusseis{toores}

\prob{Kolm vedelikku}
Reastada ained tiheduse järgi. Põhjendada tehtud valikut. NB! Ära maitse vedelikke, võivad olla ohtlikud.

\hint

\solu
Teooria: Vedeliku tiheduse määramiseks kasutame areomeetrit, mille valmistame pulgakesest ja plastiliinist. Kui üleslükkejõud on areomeetri kaaluga võrdne, siis areomeeter heljub vedelikus nii, et ta ei ulatu üle veepinna. Kui teise vedeliku tihedus on suurem, siis ulatub areomeetri ots vedelikust välja. Pöhjus on selles, et ka nüüd on tarvis kompenseerida pulga kaal. Kuid kui tihedus on suurem, siis sellevõrra väiksem ruumala vedelikku tuleb välja tõrjuda (FU = V g). Kui tihedus on väiksem, siis areomeeter upub. Areomeetri idee peale tulek (2 p). Selle tööprintsiibi seletamine nii suurema kui väiksema tihedusega vedeliku korral (2p). Valemi kasutamine põhjendamisel (1 p).

Katse: Eristatavate katsetulemuste saamine (1 p). Katsetulemuste registreerimine (klaasi number, kas ujub, heljub või upub) (1 p). Kolm ainet õiges järjekorras (2 p), kaks

ainet õiges järjekorras (1 p).

Järeldus: Katse tulemuse kokkuvõte, hinnang täpsusele või meetodile jne (1 p).
\probend
